\documentclass[11pt,a4paper]{article}
% Préambule commun aux documents de PythonIPM/documentation.
% Compilation : tectonic <fichier>.tex
% tectonic compile avec XeTeX, qui lit l'UTF-8 nativement : surtout PAS inputenc ni fontenc,
% faits pour pdfLaTeX. Ce sont eux qui décodaient mal les caractères multi-octets — « cœur »
% ressortait en « cur », et les tirets cadratins en caractères parasites.
\usepackage{lmodern}
\usepackage[french]{babel}
\usepackage{microtype}
\usepackage[a4paper,margin=2.4cm]{geometry}
\usepackage{amsmath,amssymb}
\usepackage{booktabs}
\usepackage{longtable}
\usepackage{array}
\usepackage{ragged2e}
\usepackage{xcolor}
\usepackage{tcolorbox}
\usepackage{enumitem}
\usepackage{fancyhdr}
\usepackage[hidelinks]{hyperref}

\definecolor{encadre}{HTML}{F4F6F8}
\definecolor{bordure}{HTML}{4A6FA5}
\definecolor{alerte}{HTML}{B03A2E}

% Encadré « à retenir » : la lecture non technique du passage qui précède
\newtcolorbox{retenir}[1][]{
  colback=encadre, colframe=bordure, boxrule=0.6pt, arc=2pt,
  left=8pt, right=8pt, top=6pt, bottom=6pt,
  fonttitle=\bfseries, title={#1}
}

% Encadré d'avertissement : un point ouvert, un écart assumé
\newtcolorbox{attention}[1][]{
  colback=white, colframe=alerte, boxrule=0.6pt, arc=2pt,
  left=8pt, right=8pt, top=6pt, bottom=6pt,
  % le bandeau de titre a déjà la couleur d'alerte en fond : un titre `coltitle=alerte`
  % écrirait du rouge sur du rouge. On garde le blanc par défaut, comme l'encadré « retenir ».
  fonttitle=\bfseries, title={#1}
}

\newcommand{\col}[1]{\texttt{#1}}
% \ensuremath et non \textsubscript : \ipm est employé aussi bien dans le texte que dans les
% formules, et « M\textsubscript{0} » placé en mode mathématique est une erreur — silencieuse en
% PDF, mais qui fait échouer la conversion vers Word.
\newcommand{\ipm}{\ensuremath{M_{0}}}

\setlist[itemize]{itemsep=2pt, topsep=4pt}
\setlist[description]{style=nextline, leftmargin=1.2cm, itemsep=5pt}

% La source du document apparaît en tête de chaque page. Les documents RGPH la redéfinissent
% par \renewcommand{\source}{...} AVANT \begin{document}.
\newcommand{\source}{EHCVM 2021}

\pagestyle{fancy}
\fancyhf{}
\fancyhead[L]{\small\itshape IPM Côte d'Ivoire --- \source}
\fancyhead[R]{\small\thepage}
\renewcommand{\headrulewidth}{0.3pt}

\renewcommand{\arraystretch}{1.15}

\renewcommand{\source}{RGPH 2021}

\title{\bfseries Lecture des sorties\\[4pt]
\large Ce que contiennent les classeurs de l'IPM, et comment les lire\\[2pt]
\large RGPH 2021}
\author{Pipeline \col{PythonIPM} --- IPM Côte d'Ivoire}
\date{Septembre 2026}

\begin{document}
\maketitle

\begin{abstract}
\noindent Ce document accompagne les classeurs produits par l'étape 05 du pipeline sur le
recensement. Il décrit ce qu'on y trouve, ce qui diffère des sorties de l'EHCVM, et propose une
lecture commentée des chiffres --- national, territorial, par indicateur --- avec les précautions
qui s'attachent à chacun. Les définitions colonne par colonne sont communes aux deux sources et
figurent dans le \emph{Dictionnaire des sorties} de l'EHCVM ; on n'y revient ici que là où le
recensement change quelque chose.
\end{abstract}

\begin{retenir}[Où trouver le chiffre que je cherche]
\begin{center}
\small
\begin{tabular}{@{}ll@{}}
\textbf{La question} & \textbf{Où regarder} \\[2pt]
L'IPM de la Côte d'Ivoire ? & \col{indices\_ipm\_ci\_rgph2021.xlsx}, feuille \textsf{Ensemble} \\
L'IPM d'une sous-préfecture ? & feuille \textsf{Sous-préfecture ou commune} \\
Quel territoire pèse le plus ? & \col{contribution\_M0} de la feuille correspondante \\
Quel indicateur pèse le plus ? & feuille \textsf{Contributions} \\
Le chiffre est-il publiable ? & \col{M0\_cv} --- au-delà de 0,20, s'abstenir \\
Comparer au chiffre de l'EHCVM ? & \col{indices\_ipm\_harmonise\_*} des deux sources \\
\end{tabular}
\end{center}
\end{retenir}

\tableofcontents
\bigskip

\section{Ce que le pipeline a produit}

Trois classeurs, de structure identique, qui ne diffèrent que par le périmètre d'indicateurs.

\begin{center}
\small
\begin{tabular}{@{}l r r >{\RaggedRight}p{6.0cm}@{}}
\toprule
\textbf{Classeur} & \textbf{Ind.} & \textbf{Dim.} & \textbf{Ce qu'il mesure} \\
\midrule
\col{indices\_ipm\_ci\_rgph2021} & 13 & 4 & l'\textbf{IPM national} appliqué au recensement --- le
  chiffre de référence \\
\col{indices\_ipm\_international\_rgph2021} & 11 & 3 & le périmètre comparable au niveau
  international : l'Emploi est retiré \\
\col{indices\_ipm\_harmonise\_rgph2021} & 12 & 3 & le \textbf{socle commun} avec l'EHCVM : la
  seule variante comparable d'une source à l'autre \\
\bottomrule
\end{tabular}
\end{center}

À côté des classeurs, chaque variante est aussi écrite en \col{.dta} et \col{.csv}, accompagnée de
sa table \col{contributions\_\dots}. Les tables de travail --- \col{preconstruction\_\dots},
\col{matrice\_situationnelle\_\dots}, \col{matrice\_privations\_ponderees\_\dots},
\col{matrice\_privations\_censuree\_\dots} --- ne sont écrites qu'en \col{.dta} : à 5,6 millions de
lignes, leur doublon texte pèserait plusieurs gigaoctets que personne ne lit.

\begin{retenir}[Un suffixe par source]
Tous les fichiers portent le suffixe de leur source : \col{\_rgph2021} ou \col{\_ehcvm2021}. Les
deux séries cohabitent dans le même dossier \col{sorties/} sans se recouvrir. Un fichier sans
suffixe provient d'une exécution antérieure à l'extension du pipeline au recensement.
\end{retenir}

\section{Huit feuilles, et non neuf}

\begin{center}
\begin{tabular}{@{}llr@{}}
\toprule
\textbf{Feuille} & \textbf{Une ligne =} & \textbf{Lignes} \\
\midrule
\textsf{Ensemble} & la Côte d'Ivoire entière & 1 \\
\textsf{Contributions} & un indicateur & 13 \\
\textsf{Ponderations} & un indicateur : sa dimension et son poids $w_j$ & 13 \\
\textsf{Milieu de résidence} & urbain, semi-urbain ou rural & 3 \\
\textsf{Région} & une des 33 régions & 33 \\
\textsf{Département} & un des 111 départements & 111 \\
\textsf{Sous-préfecture ou commune} & une des 519 sous-préfectures & 519 \\
\textsf{Taille du ménage} & une classe de taille & 5 \\
\bottomrule
\end{tabular}
\end{center}

Trois différences avec les sorties de l'EHCVM.

\begin{description}
  \item[Pas de feuille \textsf{Sexe du chef de ménage}.] Le sexe n'est pas dans l'extrait de
    recensement transmis. La désagrégation est simplement absente ; le pipeline la signale dans son
    journal plutôt que de produire une feuille vide.

  \item[Trois milieux au lieu de deux.] Le recensement distingue l'\textbf{urbain}, le
    \textbf{semi-urbain} et le \textbf{rural}, là où l'enquête n'en code que deux. Le semi-urbain
    n'a pas d'équivalent dans les sorties EHCVM : ne pas additionner les deux séries.

  \item[Un territoire plus fin.] 111 départements et 519 sous-préfectures ou communes, contre 108 et
    442 dans l'enquête. L'écart tient au découpage administratif retenu et au fait que le
    recensement atteint des localités qu'un échantillon ne touche pas.
\end{description}

\section{Trois colonnes qui ne se lisent pas comme dans l'enquête}

Les colonnes sont les mêmes et les formules sont les mêmes. Trois d'entre elles changent de sens
parce que la source change de nature.

\begin{description}
  \item[\col{menages}] n'est plus un effectif d'échantillon mais un \textbf{effectif réel}. La
    feuille \textsf{Ensemble} affiche 5\,616\,487 : ce sont tous les ménages du pays, pas un
    sous-ensemble tiré au sort. La sous-préfecture médiane en compte 5\,224, la plus petite 597.
    Cette colonne ne sert donc plus à juger de la robustesse d'un chiffre --- pour cela, voir
    \col{grappes} et \col{M0\_cv}.

  \item[\col{grappes}] compte les \textbf{zones de dénombrement}, l'unité de collecte du
    recensement, au nombre de 28\,329. Elle ne désigne plus une unité primaire de sondage, puisqu'il
    n'y a pas de tirage. Elle conserve son rôle dans le calcul de la variance : c'est à cette maille
    que les erreurs de dénombrement se corrèlent.

  \item[\col{population}] vaut 29\,276\,660 sur la feuille \textsf{Ensemble}. C'est le résultat de
    l'application du taux net de couverture post-censitaire (\col{EW}) aux 27\,984\,405 individus
    recensés. Le chiffre publié du RGPH 2021 est 29\,389\,150 : l'écart est de 0,4~\%.
\end{description}

\begin{attention}[Comment lire un intervalle de confiance sur un recensement]
Les colonnes \col{M0\_ic\_bas}, \col{M0\_ic\_haut}, \col{M0\_erreur\_type} et \col{M0\_cv} sont
calculées exactement comme sur l'enquête. Mais leur sens change : elles ne mesurent
\textbf{pas une incertitude d'échantillonnage} --- il n'y a pas d'échantillon --- mais la
\textbf{dispersion du phénomène entre zones de dénombrement}.

Elles restent utiles, et pour une raison précise : elles disent si une sous-préfecture au chiffre
extrême l'est réellement, ou si son chiffre est porté par deux ou trois zones atypiques. Elles sont,
par construction, bien plus étroites que celles de l'EHCVM --- l'IPM national est donné à
[0,1645 ; 0,1684] contre [0,2701 ; 0,3026] sur l'enquête. \textbf{Ne pas en conclure que le
recensement est « douze fois plus précis »} : les deux grandeurs ne mesurent pas la même chose.
\end{attention}

\section{Lecture du chiffre national}

\begin{center}
\begin{tabular}{@{}lrrr@{}}
\toprule
& $H$ & $A$ & \ipm \\
\midrule
IPM national (13 indicateurs, 4 dimensions) & \textbf{38,8~\%} & \textbf{0,4287} &
  \textbf{0,1664} \\
\quad IC 95~\% & [38,4 ; 39,3] & [0,4281 ; 0,4294] & [0,1645 ; 0,1684] \\
\bottomrule
\end{tabular}
\end{center}

\begin{retenir}[Les trois phrases à retenir]
\begin{itemize}
  \item \textbf{38,8~\%} de la population ivoirienne --- 11,4 millions de personnes --- vit dans un
        ménage qui cumule au moins un tiers des privations pondérées.
  \item Ces personnes subissent en moyenne \textbf{42,9~\%} des privations pondérées possibles.
  \item L'indice qui combine les deux vaut \textbf{0,1664}.
\end{itemize}
\end{retenir}

À côté du seuil, deux indicateurs complémentaires : \textbf{20,7~\%} de la population est
\emph{vulnérable} --- un score entre 0,2 et 1/3, donc juste sous le seuil --- et \textbf{7,8~\%} est
en \emph{pauvreté sévère}, avec un score d'au moins 0,5.

\begin{attention}[Ce chiffre ne se compare pas à l'IPM de l'EHCVM]
L'IPM national calculé sur l'enquête vaut 0,2863, celui-ci 0,1664. L'écart n'est pas une
contradiction entre deux mesures : les deux indices ne portent pas sur les mêmes indicateurs.

Le recensement ne pose aucune question de santé. Sa dimension Santé ne tient qu'à la mortalité
--- une privation qui ne touche que 1,1~\% des ménages --- et cette dimension porte pourtant
\textbf{un quart du poids}. Un ménage du recensement ne peut donc pratiquement pas dépasser un
score de 0,75, et beaucoup de ménages que l'enquête classerait pauvres restent sous le seuil.

Pour comparer les deux sources, il faut passer par le classeur \col{indices\_ipm\_harmonise\_*}
(section~\ref{sec:harmonise}).
\end{attention}

\section{Lecture territoriale}

\subsection{Le milieu de résidence}

\begin{center}
\begin{tabular}{@{}lrrrrrr@{}}
\toprule
\textbf{Milieu} & \textbf{Part pop.} & $H$ & \ipm & \textbf{Contrib.} & \textbf{Vulnér.} &
  \textbf{Sévère} \\
\midrule
Urbain & 44,8~\% & 5,4~\% & 0,0212 & 5,7~\% & 17,7~\% & 0,5~\% \\
Semi-urbain & 9,4~\% & 41,0~\% & 0,1705 & 9,7~\% & \textbf{34,2~\%} & 6,4~\% \\
Rural & 45,8~\% & 71,1~\% & 0,3075 & \textbf{84,7~\%} & 20,9~\% & 15,3~\% \\
\bottomrule
\end{tabular}
\end{center}

C'est le résultat le plus net des sorties. Le rural et l'urbain pèsent le même poids démographique
--- 45,8~\% contre 44,8~\% --- mais le rural porte \textbf{84,7~\%} de la pauvreté
multidimensionnelle du pays, contre 5,7~\% pour l'urbain. Le rapport des \ipm{} est de un à quinze.

Le semi-urbain mérite une lecture à part. Son \ipm{} est intermédiaire, mais c'est le milieu où la
\textbf{vulnérabilité est la plus forte} : 34,2~\% de sa population a un score entre 0,2 et 1/3,
contre 17,7~\% en urbain et 20,9~\% en rural. Autrement dit, une part inhabituellement large de sa
population est juste sous le seuil --- une population que peu de choses sépareraient du basculement.

\subsection{Les régions}

\begin{center}
\small
\begin{tabular}{@{}lrrrr@{}}
\toprule
\textbf{Région} & \textbf{Part pop.} & $H$ & \ipm & \textbf{Contribution} \\
\midrule
\multicolumn{5}{@{}l}{\itshape Les cinq plus faibles} \\
District autonome d'Abidjan & 21,5~\% & 1,8~\% & 0,0069 & 0,9~\% \\
District autonome de Yamoussoukro & 1,4~\% & 18,7~\% & 0,0755 & 0,6~\% \\
Sud-Comoé & 2,7~\% & 28,1~\% & 0,1181 & 1,9~\% \\
Gbêkê & 4,6~\% & 30,2~\% & 0,1279 & 3,5~\% \\
La Mé & 2,5~\% & 32,2~\% & 0,1324 & 2,0~\% \\
\addlinespace
\multicolumn{5}{@{}l}{\itshape Les cinq plus élevées} \\
Gbôklé & 1,6~\% & 65,1~\% & 0,2804 & 2,6~\% \\
Béré & 1,7~\% & 66,9~\% & 0,2843 & 2,9~\% \\
Bafing & 0,9~\% & 69,5~\% & 0,3139 & 1,7~\% \\
Folon & 0,5~\% & 76,2~\% & 0,3275 & 1,0~\% \\
Bounkani & 1,5~\% & 76,5~\% & 0,3654 & 3,2~\% \\
\bottomrule
\end{tabular}
\end{center}

L'écart va de 0,0069 à 0,3654 : un rapport de \textbf{un à cinquante-trois}. Abidjan concentre
21,5~\% de la population du pays et moins de 1~\% de sa pauvreté multidimensionnelle.

\begin{retenir}[Ne pas confondre \ipm{} et contribution]
Bounkani a l'\ipm{} le plus élevé du pays (0,3654) mais ne contribue qu'à 3,2~\% de l'indice
national, parce qu'elle ne pèse que 1,5~\% de la population. \ipm{} dit \emph{où la pauvreté est la
plus intense} ; \col{contribution\_M0} dit \emph{où se trouve le plus de pauvreté}. La première
oriente le ciblage, la seconde le volume des moyens.
\end{retenir}

\subsection{Départements et sous-préfectures}

C'est la raison d'être du calcul sur recensement. À la maille du \textbf{département}, l'\ipm{} va
d'Abidjan (0,0069) à Téhini (0,4191), en passant par Sipilou (0,3760) et Kong (0,3699). À celle de
la \textbf{sous-préfecture ou commune}, il va du Plateau (0,0015) à Santa de Ouaninou (0,4752), où
92,6~\% de la population est multidimensionnellement pauvre --- suivie de Niamoué (0,4748) et Gogo
(0,4747).

\begin{attention}[Dix-neuf sous-préfectures à ne pas publier telles quelles]
Sur les 519 modalités de la feuille \textsf{Sous-préfecture ou commune}, \textbf{19 dépassent un
coefficient de variation de 20~\%}. Elles représentent ensemble 1,55~\% de la population. Les plus
imprécises sont Plateau (CV 39~\%, 8 zones de dénombrement), Oress-Krobou (38~\%, 6 zones) et
Toukouzou (30~\%, 3 zones).

Le point commun est le petit nombre de zones de dénombrement, pas le petit nombre de ménages :
Plateau compte 1\,584 ménages, mais tous issus de 8 zones seulement. La colonne \col{grappes} est
donc le bon signal d'alerte, pas \col{menages}.

Toutes les autres mailles --- ensemble, milieu, 33 régions, 111 départements, classes de taille ---
sont sous le seuil de 20~\%, sans exception.
\end{attention}

\subsection{La taille du ménage}

\begin{center}
\begin{tabular}{@{}lrrrr@{}}
\toprule
\textbf{Taille} & \textbf{Part pop.} & $H$ & \ipm & \textbf{Contribution} \\
\midrule
1-2 personnes & 8,4~\% & 20,9~\% & 0,0825 & 4,2~\% \\
3-4 & 17,8~\% & 28,8~\% & 0,1187 & 12,7~\% \\
5-6 & 22,3~\% & 34,9~\% & 0,1491 & 19,9~\% \\
7-9 & 22,2~\% & 40,5~\% & 0,1747 & 23,3~\% \\
10 et plus & 29,3~\% & 51,8~\% & 0,2266 & \textbf{39,9~\%} \\
\bottomrule
\end{tabular}
\end{center}

\begin{attention}[Une pente en partie mécanique]
La relation est monotone, mais une partie en est un artefact de définition. Plusieurs indicateurs
sont formulés en « \emph{au moins un} membre » --- au moins un enfant non scolarisé, au moins un
chômeur, au moins un membre non alphabétisé. Un ménage de dix personnes a mécaniquement plus de
chances d'en compter un qu'un ménage de deux, à conditions de vie identiques.

Cette feuille se lit donc comme une description, non comme une causalité : elle ne dit pas que
« les grands ménages sont plus pauvres », mais que « la pauvreté mesurée par cet indice se
concentre dans les grands ménages ».
\end{attention}

\section{Lecture de la feuille \textsf{Contributions}}

Cette feuille répond à une seule question : \emph{quel indicateur fait l'indice ?} Sa colonne de
contribution somme à 1.

\begin{center}
\small
\begin{tabular}{@{}llrrr@{}}
\toprule
\textbf{Dimension} & \textbf{Indicateur} & $w_j$ & \textbf{Priv. censurée} &
  \textbf{Contribution} \\
\midrule
Emploi & Emploi de subsistance & 0,1250 & 35,7~\% & \textbf{26,8~\%} \\
Éducation & Alphabétisation & 0,0625 & 36,3~\% & 13,6~\% \\
Éducation & Année de scolarité & 0,0625 & 35,6~\% & 13,4~\% \\
Cond. de vie & Énergie de cuisson & 0,0417 & 37,8~\% & 9,5~\% \\
Éducation & Fréquentation scolaire & 0,0625 & 23,2~\% & 8,7~\% \\
Cond. de vie & Toilette & 0,0417 & 29,4~\% & 7,4~\% \\
Éducation & État civil & 0,0625 & 13,1~\% & 4,9~\% \\
Cond. de vie & Logement & 0,0417 & 18,9~\% & 4,7~\% \\
Cond. de vie & Eau potable & 0,0417 & 11,2~\% & 2,8~\% \\
Cond. de vie & Électricité & 0,0417 & 9,9~\% & 2,5~\% \\
Cond. de vie & Biens d'équipement & 0,0417 & 9,3~\% & 2,3~\% \\
Emploi & Chômage & 0,1250 & 2,4~\% & 1,8~\% \\
Santé & Mortalité & 0,2500 & 1,0~\% & \textbf{1,6~\%} \\
\bottomrule
\end{tabular}
\end{center}

\begin{center}
\begin{tabular}{@{}lrr@{}}
\toprule
\textbf{Dimension} & \textbf{Poids nominal} & \textbf{Contribution réelle} \\
\midrule
Éducation & 25~\% & 40,6~\% \\
Conditions de vie & 25~\% & 29,2~\% \\
Emploi & 25~\% & 28,6~\% \\
Santé & 25~\% & \textbf{1,6~\%} \\
\bottomrule
\end{tabular}
\end{center}

\begin{retenir}[Ce que « privation censurée » veut dire]
\col{taux\_privation\_censure} n'est pas le taux de privation brut : c'est la privation
\textbf{des seuls ménages pauvres}. C'est pour cette raison que les deux ne coïncident pas ---
l'énergie de cuisson prive 62,7~\% de la population, mais 37,8~\% en privation censurée. La
contribution se lit sur cette seconde grandeur : c'est elle qui décompose \ipm{}.
\end{retenir}

Deux lectures s'imposent devant cette feuille.

\begin{attention}[La dimension Santé est presque inerte]
Un poids nominal de 25~\%, une contribution réelle de 1,6~\%. La dimension ne tient qu'à la
mortalité : le recensement ne pose aucune autre question de santé, et la mortalité n'y est relevée
que sur douze mois quand l'IPM mondial en retient soixante. S'y ajoute que l'âge au décès est
inconnu pour 27,1~\% des décès enregistrés.

Conséquence pratique : l'IPM RGPH se comporte, pour l'essentiel, comme un indice à trois dimensions
dont la quatrième plafonne le score à 0,75. C'est la principale réserve à formuler devant le
chiffre de 0,1664.
\end{attention}

\begin{attention}[L'emploi de subsistance porte plus du quart de l'indice]
26,8~\% de \ipm{} sur un seul indicateur --- et c'est celui dont la mesure est la plus approximative.
Le recensement ne demande pas si le chef de ménage travaille « sur son propre champ » : on retient
donc le chef occupé dans une branche agricole, à son compte ou comme aide familial.

Deux effets se conjuguent pour lui donner ce poids : la dimension Emploi ne compte que deux
indicateurs, chacun à 0,125 ; et le chômage ne prive que 2,4~\% des ménages pauvres, laissant
l'emploi de subsistance porter la dimension presque seul.

Si l'approximation devait être resserrée, c'est plus du quart de l'indice qui se déplacerait. C'est
le premier point à réexaminer avant publication.
\end{attention}

\section{Comparer le recensement et l'enquête}
\label{sec:harmonise}

Les deux IPM nationaux ne se comparent pas. Le classeur \col{indices\_ipm\_harmonise\_*} existe pour
cela : il restreint le calcul au \textbf{socle commun} --- les douze indicateurs que les deux
sources savent mesurer --- en trois dimensions à 1/3. La dimension Santé disparaît : l'enquête et le
recensement n'ont pas un seul indicateur de santé en commun.

\begin{center}
\begin{tabular}{@{}lrrr@{}}
\toprule
& $H$ & $A$ & \ipm \\
\midrule
EHCVM 2021 & 53,5~\% & 0,5113 & \textbf{0,2737} \\
RGPH 2021 & 52,0~\% & 0,5157 & \textbf{0,2680} \\
\bottomrule
\end{tabular}
\end{center}

\begin{retenir}[Le meilleur contrôle de cohérence disponible]
À indicateurs identiques, une enquête par sondage de 12\,965 ménages et un recensement exhaustif de
5,6 millions de ménages, collectés la même année par deux dispositifs entièrement différents,
donnent \textbf{0,2737} et \textbf{0,2680} --- soit \textbf{2~\% d'écart}. L'incidence diffère de
1,5 point, l'intensité de 0,004.
\end{retenir}

La décomposition confirme l'accord indicateur par indicateur.

\begin{center}
\small
\begin{tabular}{@{}lrr@{}}
\toprule
\textbf{Contribution à \ipm} & \textbf{EHCVM} & \textbf{RGPH} \\
\midrule
Emploi de subsistance & 23,3~\% & 26,6~\% \\
Année de scolarité & 14,2~\% & 14,1~\% \\
Alphabétisation & 13,6~\% & 14,5~\% \\
Énergie de cuisson & 10,1~\% & 10,0~\% \\
Toilette & 9,6~\% & 7,5~\% \\
Fréquentation scolaire & 8,2~\% & 8,7~\% \\
État civil & 7,2~\% & 4,5~\% \\
Eau potable & 3,5~\% & 2,7~\% \\
Biens d'équipement & 3,3~\% & 2,5~\% \\
Logement & 3,2~\% & 4,6~\% \\
Chômage & 2,0~\% & 2,1~\% \\
Électricité & 1,8~\% & 2,3~\% \\
\bottomrule
\end{tabular}
\end{center}

Les deux seuls écarts notables sont ceux qu'on attend. La \textbf{toilette} pèse moins au
recensement (7,5~\% contre 9,6~\%) parce qu'il ne mesure pas le partage des sanitaires, second
critère de la définition PNUD. L'\textbf{état civil} pèse moins également (4,5~\% contre 7,2~\%) :
la possession d'un acte de naissance est mieux renseignée au recensement.

\begin{attention}[Ce que la variante harmonisée ne dit pas]
Elle valide la \emph{cohérence} des deux sources, pas l'IPM du pays. L'indice national de référence
reste celui de l'EHCVM, à seize indicateurs ($\ipm = 0{,}2863$). La variante harmonisée sert au
contrôle et à la comparaison entre sources ; elle n'a pas vocation à être publiée comme l'indice
de la Côte d'Ivoire.
\end{attention}

\section{Six précautions avant de citer un chiffre}

\begin{enumerate}
  \item \textbf{Ne pas comparer l'IPM RGPH (0,1664) à l'IPM EHCVM (0,2863).} Passer par la variante
        harmonisée. La différence de niveau vient de la dimension Santé, pas d'un désaccord entre
        les deux sources.

  \item \textbf{Vérifier \col{M0\_cv} avant de citer une sous-préfecture.} Dix-neuf des 519 sont
        au-delà de 20~\%. Le signal d'alerte est la colonne \col{grappes}, pas \col{menages}.

  \item \textbf{Ne pas lire les intervalles de confiance comme une précision d'échantillonnage.}
        Ils mesurent une dispersion entre zones de dénombrement. Leur étroitesse ne rend pas le
        recensement « plus précis » que l'enquête.

  \item \textbf{Distinguer \ipm{} et \col{contribution\_M0}.} Le premier localise l'intensité, le
        second le volume. Bounkani a le premier, le rural le second.

  \item \textbf{Ne pas additionner les milieux des deux sources.} Le semi-urbain du recensement n'a
        pas d'équivalent dans l'enquête.

  \item \textbf{Signaler les trois approximations} quand le chiffre est publié : l'année de
        scolarité mesurée par le diplôme, l'emploi de subsistance approché par la branche et le
        statut, l'eau et la toilette amputées de leur second critère. Elles sont détaillées dans la
        \emph{Méthodologie de calcul sur le RGPH 2021}.
\end{enumerate}

\section{Les autres sorties}

\begin{description}
  \item[\col{vecteur\_z\_rgph2021}] les treize indicateurs avec, pour chacun, sa dimension, l'énoncé
    exact appliqué, la source de la définition (proposition nationale ou application PNUD) et la
    valeur du seuil. C'est le document à consulter pour savoir ce qu'un indicateur mesure
    précisément.

  \item[\col{vecteur\_w\_rgph2021}] le poids $w_j$ de chaque indicateur et le poids de sa dimension.
    Somme à 1.

  \item[\col{matrice\_situationnelle\_rgph2021}] la matrice de privation $g^0$ : une ligne par
    ménage, une colonne 0/1 par indicateur, plus la pondération et les variables de désagrégation.
    C'est le fichier à reprendre pour toute analyse qui sortirait du cadre des classeurs --- un
    croisement inédit, une carte, un profil de privations.

  \item[\col{matrice\_privations\_ponderees\_rgph2021} et \col{\dots\_censuree\_\dots}] les états
    intermédiaires du calcul : $g^0$ pondérée et score $c_i$ pour la première, matrice censurée au
    seuil $k$ et statuts pauvre / vulnérable / sévère pour la seconde. Elles servent au contrôle et
    à la reproduction pas à pas.

  \item[\col{preconstruction\_matrice\_situationnelle\_rgph2021}] les situations brutes, avant
    application de tout seuil : pour chaque indicateur, l'effectif \emph{concerné} à côté de
    l'effectif \emph{défavorable}. C'est là qu'on va pour savoir combien de ménages n'étaient pas
    concernés par un indicateur, ou pour recalculer un taux avec un autre seuil.
\end{description}

\end{document}
